\documentclass[../../../main.tex]{subfiles}

\begin{document}

\section{Deep Neural Network}\label{Appendix:DNN}

Examples of machine learning (ML) models include logistic regression corresponding to the model $p(y|\bm{x},\bm{w})=\operatorname{Ber}(y|\sigma(\bm{w}^\top\bm{x}))$ and generalized linear models (GLMs) which generalizes to other kinds of output distributions such as Poisson. However, these models make the strong assumption that the input-output mapping is linear. A simple way of increasing the flexibility of such models is to perform a feature transformation, by replacing $\bm{x}$ with $\phi(\bm{x})$. For example, we can use a polynomial transform, given by $\phi(\bm{x})=[1,\bm{x},\bm{x}^2,\bm{x}^3,\ldots]$. This is called basis function expansion and the model becomes the following \autocite{Murphy2022}.

\begin{equation}
f(\bm{x};\bm{\theta})=\bm{W}\phi(\bm{x})+\bm{b}.
\end{equation}

The updated model still linear in the parameters $\bm{\theta}=(\bm{W},\bm{b})$. This makes model fitting easy, however, having to specify the feature transformation by hand is very limiting. A natural extension is to give the feature extractor its own parameter $\bm{\theta}_2$ giving the following where $\bm{\theta}=(\bm{\theta}_1,\bm{\theta}_2)$ and $\bm{\theta}_1=(\bm{W},\bm{b})$ \autocite{Murphy2022}.

\begin{equation}
f(\bm{x};\bm{\theta})=\bm{W}\phi(\bm{x};\bm{\theta}_2)+\bm{b}.
\end{equation}

This process can be repeated recursively to create more and more complex functions, and if we compose $L$ functions we get the following where $f_\ell(\bm{x})=f(\bm{x};\bm{\theta}_\ell)$ is the function at layer $\ell$ \autocite{Murphy2022}.

\begin{equation}
f(\bm{x};\bm{\theta}) = f_L(f_{L-1}(\cdots (f_1(\bm{x})) \cdots)).
\label{Equation:MLP}
\end{equation}

This is the key idea behind deep neural networks (DNNs). The term ``DNN'' actually encompasses a larger family of models in which we compose differentiable functions into any kind of directed acyclic graph (DAG) mapping input to output. \autoref{Equation:MLP} is the simplest example where the DAG is a chain. This is known as a feedforward neural network (FFNN) or multilayer perceptron (MLP) and assumes that the input is a fixed-dimensional vector such as $\bm{x}\in\mathbb{R}^D$ \autocite{Murphy2022}.

% Multilayer perceptron
\begin{figure}[!htbp]
\centering
\includegraphics[width=0.7\linewidth]{appendix/methods/pipelines/assets/dnn/dnn.png}
\caption[Multilayer perceptron]{A multilayer perceptron (MLP) \autocite{BrunelloN2021}.}
\end{figure}

Historically, the perceptron is a deterministic version of logistic regression with a mapping of the following form where $H(\cdot)$ is the Heaviside step function, also known as a linear threshold function \autocite{Murphy2022}.

\begin{equation}
f(\bm{x};\bm{\theta})
= \bm{1}(\bm{w}^\top\bm{x}+b \geq 0)
= H(\bm{w}^\top\bm{x}+b).
\end{equation}

In this form, the MLP is a stack of perceptrons, each of which involves the non-differentiable Heaviside function. This makes models difficult to train, which is why it has never been widely used. However, we can replace the Heaviside function $H:\mathbb{R}\to\{0,1\}$ with a differentiable activation function $\varphi:\mathbb{R}\to\mathbb{R}$. More precisely, we define the hidden units $\bm{z}_l$ at each layer $l$ to be a linear transformation of the hidden units at the previous layer passed elementwise through this activation function \autocite{Murphy2022}.

\begin{equation}
\bm{z}_l=f_l(\bm{z}_{l-1})=\varphi_l(\bm{b}_l+\bm{W}_l\bm{z}_{l-1}).
\end{equation}

The quantity that is passed to the activation function is called the pre-activations \autocite{Murphy2022}.

\begin{equation}
\bm{a}_l=\bm{b}_l+\bm{W}_l\bm{z}_{l-1}.
\end{equation}

If we now compute $L$ of these functions together, as in \autoref{Equation:MLP}, then we can compute the gradient of the output with respect to the parameters in each layer using the chain rule, also known as backpropagation. We can then pass the gradient to an optimizer, and thus minimize some training objective. The term ``MLP'' almost always refers to this differentiable form of the model, rather than the historical version with non-differentiable linear threshold units \autocite{Murphy2022}.

We are free to use any kind of differentiable activation function at each layer. However, if we use a linear activation function, then the whole model reduces to a regular linear model. In the early days of neural networks, a common choice of activation function was to use a sigmoid (logistic) function, which can be seen as a smooth approximation of the Heaviside function \autocite{Murphy2022}.

\begin{equation}
\sigma(a) = \frac{1}{1+e^{-a}}.
\end{equation}

However, the sigmoid function saturates at 1 for large positive inputs, and at 0 for large negative inputs. Another common choice is the tanh function, which has a similar shape, but saturates at -1 and 1. In the saturated regimes, the gradient of the output with respect to the input will be close to zero, so any gradient signal from higher layers will not be able to propagate back to earlier layers. This is called the vanishing gradient problem, and it makes it hard to train the model using gradient descent \autocite{Murphy2022}.

One of the keys to train very deep models is to use non-saturating activation functions. The most common of these is rectified linear unit (ReLU). The ReLU functions simply turns off negative inputs, and passes positive inputs unchanged \autocite{Murphy2022}.

\begin{equation}
\operatorname{ReLU}(a)
= \max(a,0)
= a\bm{1}(a>0).
\end{equation}

It can be shown that an MLP with one hidden layer is a universal function approximator, meaning it can model any suitably smooth function, given enough hidden units, to any desired level of accuracy. Intuitively, the reason for this is that each hidden unit can specify a half plane, and a sufficiently large combination of these can divide any region of space, to which any response can be associated \autocite{Murphy2022}.

% Vornoi diagram
\begin{figure}[!htbp]
\centering
\includegraphics[width=0.4\linewidth]{appendix/methods/pipelines/assets/dnn/voronoi.png}
\caption[Voronoi diagram]{Voronoi diagram visualization of the universal function approximation where each hidden unit specifies a half plane \autocite{Jahobr2017}.}
\end{figure}

However, various arguments, both experimental and theoretical, have shown that deep networks work better than shallow ones. The reason is that later layers can leverage the features that are learned by earlier layers. That is, the function is defined in a compositional or hierarchical way. Although we could fit any dataset with a single hidden layer model, intuitively, it is easier if the model learns to detect the separating characteristics of the dataset using the hidden units in an early layer, and then uses these features to define a simple linear classifier in a later layer \autocite{Murphy2022}.

\subsection{Convolutional Neural Network}

The convolutional neural network (CNN) is used to apply neural networks for images. In this case, we replace matrix multiplication with a convolution operation. The basic idea is to divide the input into overlapping 2D image patches, and to compare each patch with a set of small weight matrices, or filters, which represents part of an object \autocite{Murphy2022}.

This can be thought of as a form of template matching. We learn these templates from data. Since the templates are small (often just $3 \times 3$ or $5 \times 5$), the number of parameters needed for learning is small. Additionally, since we use convolution to do template matching, instead of matrix multiplication, the model will be translationally invariant. This is useful for tasks such as image classification, where the goal is to classify if an object is present regardless of its location \autocite{Murphy2022}.

The 1D convolution between two function $f,g:\mathbb{R}^D\to\mathbb{R}$ is defined as follows \autocite{Murphy2022}.

\begin{equation}
[f \circledast g](z) = \int_{\mathbb{R}^D}f(u)g(z-u)\mathrm{d}u.
\label{Equation:ConvolutionIntegral}
\end{equation}

Now suppose we replace the functions with finite-length vectors, which we can think of as functions defined on a finite set of points. Let $f$ be evaluated at the points $\{-L,-L+1,\ldots,0,1,\ldots,L\}$ to yield the weight vector (also called a filter or kernel) $w_{-L}=f(-L)$ up to $w_L=f(L)$. Now let $g$ be evaluated at the points $\{-N,\ldots,N\}$ to yield the feature vector $x_{-N}=g(-N)$ up to $x_N=g(N)$. Then \autoref{Equation:ConvolutionIntegral} becomes the following. We see that we flip the weight vector $\bm{w}$, and then drag it over the the $\bm{x}$ vector, summing up the local windows at each point \autocite{Murphy2022}.

\begin{equation}
[\bm{w}\circledast\bm{x}](i) = w_{-L}x_{i+L}+\cdots+w_{-1}x_{i+1}+w_0x_i+w_1x_{i-1}+\cdots+w_Lx_{i-L}.
\end{equation}

When we do not flip $\bm{w}$ first, then we get the cross correlation \autocite{Murphy2022}.

\begin{equation}
[\bm{w}\circledast\bm{x}](i) = w_{-L}x_{i-L}+\cdots+w_{-1}x_{i-1}+w_0x_i+w_1x_{i+1}+\cdots+w_Lx_{i+L}.
\end{equation}

If the weight vector is symmetric, then cross correlation and convolution are the same. In the deep learning literature, the term ``convolution'' is often used to mean cross correlation. We can also evaluate the weights $\bm{w}$ on the domain $\{0,1,\ldots,L-1\}$ and the features $\bm{x}$ on the domain $\{0,1,\ldots,N-1\}$, to eliminate negative indices \autocite{Murphy2022}.

\begin{equation}
[\bm{w}\circledast\bm{x}](i) = \sum_{u=0}^{L-1}w_ux_{i+u}.
\label{Equation:1DConvolution}
\end{equation}

% 1D convolution
\begin{figure}[!htbp]
\centering
\includegraphics[width=0.8\linewidth]{appendix/methods/pipelines/assets/dnn/conv1d-diagram.png}
\caption[1D discrete convolution]{1D discrete convolution of a signal with a kernel.}
\end{figure}

In 2D, \autoref{Equation:1DConvolution} becomes the following where the 2D filter $\bm{W}$ has size $H \times W$ \autocite{Murphy2022}.

\begin{equation}
[\bm{W}\circledast\bm{X}](i,j) = \sum_{u=0}^{H-1}\sum_{v=0}^{W-1}w_{u,v}x_{i+u,j+v}.
\end{equation}

% 2D convolution
\begin{figure}[!htbp]
\centering
\includegraphics[width=0.6\linewidth]{appendix/methods/pipelines/assets/dnn/conv2d-diagram.png}
\caption[2D discrete convolution]{2D discrete convolution of a signal with a kernel.}
\end{figure}

We can think of 2D convolution as template matching, since the output at a point $(i,j)$ will be large if the corresponding image centered on $(i,j)$ is similar to $\bm{W}$. More generally, we can think of convolution as a form of feature detection. The resulting output $\bm{Y}=\bm{W}\circledast\bm{X}$ is therefore called a feature map \autocite{Murphy2022}.

In \autoref{Figure:2DConvolution} we see that convolving a $3 \times 3$ image with a $2 \times 2$ filter results in a $2 \times 2$ output. This is called valid convolution, since we only apply the filter to valid parts of the input, and we do not let it slide off the ends. If we want the output to have the same size as the input, we can use zero-padding, where we add a border of zeros to the image. This is called same convolution \autocite{Murphy2022}.

% Padding without strides
\begin{figure}[!htbp]
\centering
\begin{subfigure}{0.25\textwidth}
\includegraphics[width=\linewidth]{appendix/methods/pipelines/assets/dnn/padding-no-strides/frame00.png}
\end{subfigure}
\hspace{1em}
\begin{subfigure}{0.25\textwidth}
\includegraphics[width=\linewidth]{appendix/methods/pipelines/assets/dnn/padding-no-strides/frame01.png}
\end{subfigure}
\hspace{1em}
\begin{subfigure}{0.25\textwidth}
\includegraphics[width=\linewidth]{appendix/methods/pipelines/assets/dnn/padding-no-strides/frame02.png}
\end{subfigure}
\\
\begin{subfigure}{0.25\textwidth}
\includegraphics[width=\linewidth]{appendix/methods/pipelines/assets/dnn/padding-no-strides/frame03.png}
\end{subfigure}
\hspace{1em}
\begin{subfigure}{0.25\textwidth}
\includegraphics[width=\linewidth]{appendix/methods/pipelines/assets/dnn/padding-no-strides/frame04.png}
\end{subfigure}
\hspace{1em}
\begin{subfigure}{0.25\textwidth}
\includegraphics[width=\linewidth]{appendix/methods/pipelines/assets/dnn/padding-no-strides/frame05.png}
\end{subfigure}
\caption[Convolutions with padding]{The first 6 of 25 convolutions with padding \autocite{Dumoulin2016}.}
\end{figure}

Since each output pixel is generated by a weighted combination of inputs in its receptive field, neighboring outputs will be very similar in value, since their inputs are overlapping. We can reduce this redundancy and speedup computation by skipping every $s$ number of inputs. This is called strided convolution \autocite{Murphy2022}.

% Padding with strides
\begin{figure}[!htbp]
\centering
\begin{subfigure}{0.25\textwidth}
\includegraphics[width=\linewidth]{appendix/methods/pipelines/assets/dnn/padding-strides/frame00.png}
\end{subfigure}
\hspace{1em}
\begin{subfigure}{0.25\textwidth}
\includegraphics[width=\linewidth]{appendix/methods/pipelines/assets/dnn/padding-strides/frame01.png}
\end{subfigure}
\hspace{1em}
\begin{subfigure}{0.25\textwidth}
\includegraphics[width=\linewidth]{appendix/methods/pipelines/assets/dnn/padding-strides/frame02.png}
\end{subfigure}
\\
\begin{subfigure}{0.25\textwidth}
\includegraphics[width=\linewidth]{appendix/methods/pipelines/assets/dnn/padding-strides/frame03.png}
\end{subfigure}
\hspace{1em}
\begin{subfigure}{0.25\textwidth}
\includegraphics[width=\linewidth]{appendix/methods/pipelines/assets/dnn/padding-strides/frame04.png}
\end{subfigure}
\hspace{1em}
\begin{subfigure}{0.25\textwidth}
\includegraphics[width=\linewidth]{appendix/methods/pipelines/assets/dnn/padding-strides/frame05.png}
\end{subfigure}
\caption[Convolutions with padding and strides]{The first 6 in 9 convolutions with padding and strides \autocite{Dumoulin2016}.}
\end{figure}

Convolution will preserve information about the location of input features, a property known as equivariance. In some cases we want to be invariant to the location. For example, when performing image classification, we may just want to know if an object of interest is present anywhere in the image. A simple way to achieve this is called max pooling, which computes the maximum over its incoming values. An alternative is to use average pooling, which replaces the maximum with the mean. A common design pattern is to create a CNN by alternating convolutional layers with max pooling layers, followed by a final linear classification layer at the end \autocite{Murphy2022}.

% Max pooling
\begin{figure}[!htbp]
\centering
\includegraphics[width=0.6\linewidth]{appendix/methods/pipelines/assets/dnn/max-pooling.png}
\caption[Max pooling]{Max pooling \autocite{Aphex342015}.}
\end{figure}

\subsection{Bayesian Neural Network}

Deep neural networks (DNNs) are usually trained using a regularized maximum likelihood objective to find a single setting of parameters. However, large flexible models like neural networks can represent many functions, corresponding to different parameter settings, which fit the training data well, yet generalize in different ways. This phenomenon is known as underspecification \autocite{Murphy2023}.

Considering all of these different models together can lead to improved accuracy and uncertainty representation. This can be done by computing the posterior predictive distribution using Bayesian model averaging. The main challenges in applying Bayesian inference to DNNs are specifying suitable priors, and efficiently computing the posterior, which is challenging due to the large number of parameters and the large datasets \autocite{Murphy2023}.

\begin{equation}
p(y|\bm{x},\mathcal{D}) = \int p(y|\bm{x},\bm{\theta})p(\bm{\theta}|\mathcal{D})\mathrm{d}\bm{\theta} \quad \text{where} \quad p(\bm{\theta}|\mathcal{D}) \propto p(\bm{\theta})p(\mathcal{D}|\bm{\theta}).
\end{equation}

Consider the generalized MLP with $L-1$ hidden layers and a linear output as described in \autoref{Equation:MLP} \autocite{Murphy2023}.

\begin{equation}
f(\bm{x};\bm{\theta}) = \bm{W}_L(\cdots\varphi(\bm{W}_1\bm{x}+\bm{b}_1))+\bm{b}_L.
\end{equation}

The most common choice of priors is to use a factored Gaussian prior \autocite{Murphy2023}.

\begin{equation}
\bm{W}_\ell \sim \mathcal{N}(\bm{0},\alpha_\ell^2\bm{I}), \quad \bm{b}_\ell \sim \mathcal{N}(\bm{0},\beta_\ell^2\bm{I}).
\end{equation}

Initializing this model's parameters at a particular random value is like sampling a point from this prior over functions. In the limit of infinitely wide neural networks, a neural network defines a Gaussian process (see \autoref{Appendix:GP}) with a fixed kernel. Indeed, an MLP with one hidden layer, whose width goes to infinity, and which has a Gaussian prior on all the parameters, converges to a Gaussian process with a well-defined kernel. These kernels can be used in modeling non-stationary covariance structure \autocite{Murphy2023}. 

Monte Carlo dropout (MCD) is a very simple and widely used method for approximating the Bayesian predictive distribution. Usually stochastic dropout layers are added as a form of regularization during training, and are turned off at test time. However, the idea of MCD is to also perform random sampling during testing \autocite{Murphy2023}.

More precisely, we drop out each hidden unit according to a $\operatorname{Bernoulli}(p)$ distribution, and repeat this procedure $K$ times to create $K$ distinct models. We then create an equally weighted average of the predictive distribution for each of these models as shown below. One drawback of MCD is that it is slow at test time. However, this can be overcome by distilling the model's predictions into a deterministic student network \autocite{Murphy2023}.

\begin{equation}
p(y|\bm{x},\mathcal{D}) \approx \frac{1}{K}\sum_{k=1}^{K}p(y|\bm{x},\bm{\theta}^k).
\end{equation}

Another very simply approximation to the posterior is to only be Bayesian about the weights in the final layer. This is called the ``Bayesian last layer'' approximation. In more detail, let $\bm{z}=f(\bm{x};\bm{\theta})$ be the predicted outputs of the model before any optional final nonlinearity. We assume this has the form $\bm{z}=\bm{w}_L^\top\phi(\bm{x};\bm{\theta})$, where $\phi(\bm{x})$ are the features extracted by the first $L-1$ layers. Then, we can use standard techniques, such as Markov chain Monte Carlo (described by \autoref{Appendix:MCMC}), to compute $p(\bm{w}_L|\mathcal{D})=\mathcal{N}(\bm{\mu}_L,\bm{\Sigma}_L)$, given $\phi(\cdot)$ \autocite{Murphy2023}.

MCMC methods like Hamiltonian Monte Carlo (HMC) are generally considered to be the gold standard for posterior approximation, since they do not make strong assumptions about the form of the posterior. However, a significant limitation of standard MCMC procedures, including HMC, is that they require access to the full training set at each step. Stochastic gradient MCMC methods, such as Stochastic gradient Langevin dynamics (SGLD), operate instead using mini-batches of data, offering a scalable alternative \autocite{Murphy2023}.

Many conventional approximate inference methods, such as variational inference, focus on approximating the posterior $p(\bm{\theta}|\mathcal{D})$ in a local neighborhood around one of the posterior modes. While this is often not a major limitation in classical machine learning, modern deep neural networks have highly multi-modal posteriors, with parameters in different modes giving rise to very different functions \autocite{Murphy2023}.

On the other hand, the functions in a neighborhood of a single mode may make fairly similar predictions. So using a local approximation to compute the posterior predictive will underestimate uncertainty and generalize more poorly. A simple alternative method is to train multiple models, and then to approximate the posterior using an equally weighted mixture of delta functions as follows where $M$ is the number of models. This approach is called deep ensembles \autocite{Murphy2023}.

\begin{equation}
p(\bm{\theta}|\mathcal{D}) \approx \frac{1}{M}\sum_{m=1}^{M} \delta(\bm{\theta}-\hat{\bm{\theta}}_m).
\end{equation}

Once we have an approximation $q(\bm{\theta}|\mathcal{D})$ of the parameter posterior $p(\bm{\theta}|\mathcal{D})$, we can use it to approximate the posterior predictive distribution \autocite{Murphy2023}.

\begin{equation}
p(y|\bm{x},\mathcal{D})=\int p(y|\bm{x},\bm{\theta})p(\bm{\theta}|\mathcal{D})\mathrm{d}\bm{\theta}.
\end{equation}

We often approximate this integral using Monte Carlo where $\bm{\theta}^s \sim q(\bm{\theta}|\mathcal{D})$ \autocite{Murphy2023}.

\begin{equation}
p(y|\bm{x},\mathcal{D}) \approx \frac{1}{S}\sum_{s=1}^{S}p(y|\bm{x},\bm{\theta}^s).
\end{equation}

\end{document}
